\chapter{Geometry}

We will start with a definition that originally comes much later in the book.
There, space is defined as the set of all points (\cite{SMSGGeometry}, 53).
Together with the earlier undefined objects points, lines and planes, of which lines and planes were loosely defined as sets of points (\cite{SMSGGeometry}, 15), we conclude that we can build the following geometry based on $\RR^3$.


\begin{definition}[Space]\label{def:Space}
  \lean{structure:Space}
  \leanok
    Space, the set of all points, is $\RR^3$. (\cite{SMSGGeometry}, 53)
\end{definition}

The SMSG further defines all geometric items (including lines and planes) as sets of points (\cite{SMSGGeometry}, 15).
Hence, we get

\begin{definition}[Sets of Points]\label{def:GeomItems}
  \lean{structure:Line}
  \leanok
  \uses{def:Space}
    For all geometric items $G$, let $G \subset \RR^3$.
\end{definition}

%_______________________ Postulate 1 _____________________________________________

\begin{definition}[Postulate 1]\label{def:Postulate_1}
  \lean{structure:ObjectsWithCoordinate.lineBy, structure:ObjectsWithCoordinate.exists_lineBy, structure:ObjectsWithCoordinate.mem_lineBy_left, structure:ObjectsWithCoordinate.mem_lineBy_right, structure:ObjectsWithCoordinate.unique_lineBy}
  \leanok
  \uses{def:GeomItems}
  Given any two different points, there is exactly one line which contains both of them. (\cite{SMSGGeometry}, 30)

  (For $P,Q \in \RR^3$, $P \ne Q$, exists exactly one line $L \subset \RR^3$ with $P,Q \in L$.
  For all other lines $L' \subset \RR^3$ with $P,Q \in L'$ we see $L=L'$.)
\end{definition}
%__________________________________________________________________________________
  To designate the line determined by two points $P,Q \in \RR^3$ we use the notation $\geomline{PQ}$.


%_________________________ Postulate 2 _________________________________________
\begin{definition}[Postulate 2 - The Distance Postulate]\label{def:Postulate_2}
  \lean{structure:Space.dist, structure:Space.dist_self, structure:Space.dist_pos}
  \leanok
  \uses{def:GeomItems}
  To every pair of points there corresponds a unique non-negative number. (\cite{SMSGGeometry}, 34)

  If the points are equal, the number is zero.
  If the points are different, the number is positive.

  (Existence of function $\operatorname{d}: (\RR^3\times\RR^3) \to \RR_+ \cup \set{0}$.)
  %TODO, add when zero and when positive
\end{definition}
%________________________________________________________________________________

\begin{definition}[Distance]\label{def:Distance}
  \lean{structure:Space.dist}
  \leanok
  \uses{def:Postulate_2}
  The distance between two points is the positive number given by the Distance Postulate (\ref{def:Postulate_2}).
  If the points are $P,Q \in \RR^3$, then the distance is denoted by $PQ$.
  (\cite{SMSGGeometry}, 34)

  ($PQ=d(P,Q)$)
\end{definition}

%____________________________ Postulate 3 ____________________________________


\begin{definition}[Postulate 3 - The Ruler Postulate]\label{def:Postulate_3}
  \lean{structure:LineWithCoordinate, structure:ObjectsWithCoordinate.line_dist}
  \leanok
  \uses{def:GeomItems, def:Postulate_1, def:Distance, def:Postulate_2}
  The points of a line can be placed in correspondence with the real numbers in such a way that
  \begin{enumerate}
    \item to every point of the line there corresponds exactly one real number,
    \item to every real number there corresponds exactly one point of the line, and
    \item the distance between two points is the absolute value of the difference of the corresponding numbers.
  \end{enumerate}
  (\cite{SMSGGeometry}, 36)

  %TODO

  For line $L \subset \RR^3$, exists bijective function $\operatorname{coord}:L \to \RR$ (covers item 1 and 2).

  Distance between two points $P,Q$ then is $\operatorname{d}=|\operatorname{coord}(P)-\operatorname{coord}(Q)|$.
  With $|\cdot |$ being the absolute value.
\end{definition}
%___________________________________________________________________________________

  This postulate also tells us that there are infinitely many points in a line.


\begin{definition}[Coordinate System]\label{def:Coordinate}
  \lean{sturcture:LineWithCoordinate.coordinate}
  \leanok
  \uses{def:Postulate_3}
  A correspondence of the sort described in the Ruler Postulate (\ref{def:Postulate_3}) is called a coordinate system
  for the line.
  The number corresponding to a given point is called the coordinate of the point.
  (\cite{SMSGGeometry}, 37)


\end{definition}

%____________________________________ Postulate 4 _________________________________________
\begin{definition}[Postulate 4 - The Ruler Placement Postulate]\label{def:Postulate_4}
  %TODO \lean{}
  \uses{def:GeomItems, def:Coordinate, def:Postulate_3}
  Given two points $P$ and $Q$ of a line, the coordinate system can be chosen in such a way that the coordinate of $P$ is zero and the coordinate of $Q$ is positive.
  (\cite{SMSGGeometry}, 40)


  For $L \subset \RR^3$ line and $P,Q \in L$, we can choose a bijective function (as in \ref{def:Postulate_3}) $\operatorname{coord}$ so that $\operatorname{coord}(P)=0$ and $\operatorname{coord}(Q)>0$.
\end{definition}
%___________________________________________________________________________________________

\begin{definition}[Betweenness]\label{def:Betweenness}
  \lean{def:between, structure:Objects}
  \leanok
  \uses{def:Postulate_1, def:Postulate_2, def:Distance} %TODO: does it use def:Postulate_3?
  For three points $A,B,C$, $B$ is between $A$ and $C$ if
  \begin{enumerate}
    \item $A,B,C$ are distinct points on the same line, and
    \item $AB + BC =AC$.
  \end{enumerate}
  (\cite{SMSGGeometry}, 41)

  For $A,B,C \in \RR^3$, $B$ is between $A$ and $C$ if (1) $A\ne B\ne C$ and $B \in \geomline{AC}$ and (2) $AB+BC=AC$ ($\operatorname{d}(A,B)+\operatorname{d}(B,C)+\operatorname{d}(A,C)$).

  We will write $A(B)C$ for $B$ is between $A$ and $C$.
\end{definition}

\begin{theorem}[2-1.]\label{thm:2-1.}
  \lean{thm:coordinates_between_points_between}
  \leanok
  \uses{def:Postulate_1, def:Postulate_3, def:Betweenness, def:Coordinate, def:Distance}
  Let $A,B,C$ be three points of a line $L$ ($A,B,C \in L$) with coordinates $x,y,z \in \RR$.
  If $x<y<z$ then $B$ is between $A$ and $C$.
  (\cite{SMSGGeometry}, 42)

  Let $A,B,C \in \RR^3$ with $B \in \geomline{AC}$, with $\coord{A}=:x, \coord{B}=:y, \coord{C}=:z$. If $x<y<z$, then $B$ is between $A$ and $C$.

\end{theorem}

\begin{proof}
  \uses{def:Postulate_1, def:Postulate_3, def:Betweenness, def:Coordinate, def:Distance}
  \leanok
    \begin{align*}
      x<y<z &\imp
                  y-x,z-y,z-x>0 \\
            &\imp
                  y-x=\abs{y-x}=AB, z-y=\abs{z-y}=BC, z-x=\abs{z-x}=AC \\
            &\imp
                  AB+BX=(y-x)+(z-y)=z-x=AC \MP.
    \end{align*}
\end{proof}

\begin{theorem}[2-2.]\label{thm:2-2.}
  \lean{thm:exists_between}
  \leanok
  \uses{def:Postulate_1, def:Postulate_3,def:Betweenness, def:Coordinate, thm:2-1.}
  Of any three different points on the same line, one is between the other two.
  (\cite{SMSGGeometry}, 43)

  For $A,B,C \in \RR^3, A\ne B\ne C$ with $B \in \geomline{AC}$: $A(B)C, B(A)C$ or $A(C)B$.
\end{theorem}

\begin{proof}
  \uses{def:Postulate_1, def:Postulate_3,def:Betweenness, def:Coordinate, thm:2-1.}
  \leanok
  Let $A,B,C \in \RR^3, A\ne B\ne C$ and $B \in \geomline{AC}$.
  \ref{def:Postulate_3} allows for coordinates, let $\coord{A}=:x, \coord{B}=:y, \coord{C}=:z$.
  Since $A \ne B \ne C$, we know $x \ne y \ne z$.

  There are now six possibilities
  \begin{align*}
    (1)& x<y<z\MP, & (2)& x<z<y\MP, & (3)& y<x<z \MP, \\
    (4)& y<z<x \MP, & (5)& z<x<y \MP, & (6)& z<y<x \MP.
  \end{align*}
  \ref{thm:2-2.} then follows from \ref{thm:2-1.}.
\end{proof}

\begin{theorem}[2-3.]\label{thm:2-3.}
  \lean{thm:unique_between}
  \leanok
  \uses{def:Postulate_1, def:Distance, def:Betweenness, thm:2-2.}
  Of three different points on the same line, only one is between the other two.
  (\cite{SMSGGeometry}, 44)

  Of $A,B,C \in \RR^3, A \ne B \ne C$ with $B \in \geomline{AC}$ either $A(B)C, A(C)B$, or $B(A)C$.
\end{theorem}
\textit{Restatement:} If $A,B$ and $C$ are three different points on the same line, and $B$ is between $A$ and $C$, than $A$ is not between $B$ and $C$, and $C$ is not between $A$ and $B$.

\begin{proof}
  \uses{def:Postulate_1, def:Distance, def:Betweenness}
  \leanok
  If $A(B)C$, then $AB+BC=AC$ (1), if additionally $B(A)C$, then $BA+AC=AB+AC=BC$.

  ((1) $\imp AC-BC=AB$ and (2) $\imp AC-BC=-AB$)

  $\imp AB=AC-BC=-AB$.

  Since $AB\ne-AB$ only one of them can be between the other two.
  The other cases follow accordingly.
\end{proof}


\begin{definition}[Segment]\label{def:Segment}
  \uses{def:Betweenness}
  For any two points $A$ and $B$ the segment $\segment{AB}$ is the set whose points are $a$ and $B$, together with all points that are between $A$ and $B$.
  The points $A$ and $B$ are called the end-points of $\segment{AB}$.
  (\cite{SMSGGeometry},45)

  For $A,B \in \RR^3: \segment{AB}:=\set{A,B}\cup \set{C\mid A(C)B}$.
\end{definition}

\begin{definition}[Segment Lenth]\label{def:Segment_Lenth}
  \uses{def:Distance, def:Segment}
    The distance $AB$ is called the lenth of the segment $\segment{AB}$.
    (\cite{SMSGGeometry}, 45)

    %TODO restatement necessary?
\end{definition}

\begin{definition}[Ray]\label{def:Ray}
  \uses{def:Postulate_1, def:Segment, def:Betweenness} %TODO does it use thm:2-3.?
  Let $A$ and $B$ be points of a line $L$.
  The ray $\ray{AB}$ is the set which is the union of (1) $\segment{AB}$ and (2) the set of all points $C$ with $B$ is in between $A$ and $C$.
  $A$ is called the end-point of $\ray{AB}$.
  (\cite{SMSGGeometry}, 46)

  Let $L \subset \RR^3$ and $A,B \in L$: $\ray{AB}:=\segment{AB}\cup \set{C\mid A(B)C}$.
\end{definition}

\begin{definition}[Opposite Rays]\label{def:Oppo_Rays}
  \uses{def:Ray, def:Betweenness}
  If $A$ is in between $B$ and $C$, then $\ray{AB}$ and $\ray{AC}$ are called opposite rays.
  (\cite{SMSGGeometry}, 46)

  %TODO restatement necessary? if so how? maybe \ray{AB}\setminus\set{A}\cup \ray{AC}\setminus\set{A}=\emptyset ?
\end{definition}

\begin{theorem}[2-4. - The Point Plotting Theorem]\label{thm:2-4.}
  \uses{def:Distance, def:Postulate_3, def:Postulate_4, def:Ray}
  Let $\ray{AB}$ be a ray and let $x$ be a positive number.
  Then there is exactly one point $P$ of $\ray{AB}$ such that $AP=x$.
  (\cite{SMSGGeometry}, 46)

  Let $\ray{AB} \subset \RR^3$ and $x \in \RR_+$ then exists exactly one $P \in \ray{AB}$ with $AP=x$.
\end{theorem}

\begin{proof}
  \uses{def:Distance, def:Postulate_3, def:Postulate_4, def:Ray}
  By \ref{def:Postulate_4} we can choose coordinate system (bijective function) on $\geomline{AB}$ such that $\coord{A}=0$ and $\coord{B}=:r>0$.
  Let $P$ be the point with $\coord{P}=x$ then $ P \in \geomline{AB}$ and $AP=\abs{x-0}=\abs{x}=x$.
  Since $\coord{}$ bijective (\ref{def:Postulate_3}) follows for all $P' \in \ray{AB}$ with $\coord{P'}=x$ that $P=P'$.
\end{proof}

\begin{definition}[Mid-Point]\label{def:Mid-Point}
  \uses{def:Distance, def:Betweenness, def:Segment}
  A point $B$ is called mid-point of a segment $\segment{AC}$, if $B$ is between $A$ and $C$ and $AB=BC$.
  (\cite{SMSGGeometry}, 47)

  $B \in \RR^3$ is mid-point of $\segment{AC} \subset \RR^3$ if $A(B)C$ and $AB=BC$.
\end{definition}

\begin{theorem}[2-5.]\label{thm:2-5.}
  \uses{def:Betweenness, def:Segment, thm:2-4., def:Mid-Point}
  Every segment has exactly one mid-point.
  (\cite{SMSGGeometry}, 47)

  For every segment $\segment{AC} \subset \RR^3$ exists exatly one $B\in \RR^3$ with $A(B)C$ and $AB=BC$.
\end{theorem}
\begin{proof}
  \uses{def:Betweenness, def:Segment, thm:2-4., def:Mid-Point, def:Ray}
  On $\segment{AC} \subset \RR^3$ choose $B \in \segment{AC}$ so that $AB=BC$ (\ref{thm:2-4.}) (1).
  By definition of segment (\ref{def:Segment}) we know $A(B)C$, hence $AB+BC=AC$ (2).
  From (1) and (2) we conclude $2AB=AC$ or $AB=\frac{1}{2}AC$.
  Since $B \in \segment{AC}$ if follows that $B \in \ray{AC}$ and \ref{thm:2-4.} says there is only exactly one such point.

  %TODO proof feels wonky
\end{proof}

\begin{definition}[Bisection]\label{def:Bisection}
  \uses{def:Segment, def:Mid-Point, thm:2-5.}
  The mid-point of a segment is said to bisect the segment.
  Any figure whose intersection with a segment is the mid-point is said to bisect the segment.
  (\cite{SMSGGeometry}, 47)
\end{definition}

\begin{definition}[Collinear]\label{def:Collinear}
  \uses{def:GeomItems}
  A set of points is collinear if there is a line which contains all the points of the set.
  (\cite{SMSGGeometry}, 54)

  %TODO restatement necessary?
\end{definition}

\begin{definition}[Coplanar]\label{def:Coplanar}
  \uses{def:GeomItems}
  A set of points is coplanar if there is a plane which contains all the points of the set.
  (\cite{SMSGGeometry}, 54)

  %TODO restatement necessary?
\end{definition}

%______________________________________ Postulate 5 _______________________________________________________
\begin{definition}[Postulate 5]\label{def:Postulate_5}
  \uses{def:Collinear, def:Coplanar} %TODO restatement also uses Postulate 1
  \begin{itemize}
    \item Every plane contains at least three non-collinear points.
    \item Space contains at least four non-coplanar points.
  \end{itemize}
  (\cite{SMSGGeometry}, 54)

  Let $E \subset \RR^3$ be a plane then exist $A, B, C \in E$ with $A \not\in \geomline{BC}$.

  There are $A,B,C,D \in \RR^3$ with $A,B,C \in M$ so that $A \not\in \geomline{BC}$ and $D \not\in M$.
\end{definition}
%________________________________________________________________________________________________________

\begin{theorem}[3-1.]\label{thm:3-1.}
  \uses{def:Postulate_1}
  Two different lines intersect in at most one point.
  (\cite{SMSGGeometry}, 55)

  For lines $L_1,L_2 \subset \RR^3$ with $L_1 \ne L_2: \#(L_1\cap L_2)\leq 1$.
\end{theorem}
\begin{proof}
  \uses{def:Postulate_1}
  Let $L_1,L_2 \subset \RR^3$ with $L_1\ne L_2$, and further $\#(L_1\cap L_2)>1$.
  Then exist $P,Q \in L_1\cap L_2$.
  Then $P,Q \in L_1$ and $P,Q \in L_2$.
  Thus, $L_1=\geomline{PQ}=L_2$.
\end{proof}

%____________________________________________ Postulate 6 _____________________________________________
\begin{definition}[Postulate 6]\label{def:Postulate_6}
  \uses{def:Postulate_1}
  If two points lie in the same plane, then the line containing these points lies in the same plane.
  Let $E \subset \RR^3$ be a plane. If $A,B \in E$, then $\geomline{AB} \subset E$
  (\cite{SMSGGeometry}, 56)
\end{definition}
%______________________________________________________________________________________________________

\begin{theorem}[3-2.]\label{3-2.}
    \uses{def:Postulate_1, def:Postulate_6}
    If a line intersects a plane not containing it then the intersection is a point.
    (\cite{SMSGGeometry}, 56)

    Let $L \subset \RR^3$ be a line and $E \subset \RR^3$ be a plane.
    If $L \not\subset E$ and $\#(L\cap E)>0$ then $\#(L\cap E)=1$.
\end{theorem}

\begin{proof}
  \uses{def:Postulate_1, def:Postulate_6}
  Assume $\#(L\cap E)>1$, then exist points $A,B \in L\cap E$.
  Then according to \ref{def:Postulate_6} $L=\geomline{AB}\subset E$.
\end{proof}

%_________________________________ Postulate 7 ______________________________________________________
\begin{definition}[Postulate 7]\label{def:Postulate_7}
  \uses{def:Postulate_1, def:Collinear, def:Coplanar}
  Any three points lie in at least one plane, and any three non-collinear points liein exactly one plane.
  More briefly, any three points are coplanar and any three non-collinear points determine a plane.
  (\cite{SMSGGeometry}, 57)

  For $A,B,C \in \RR^3$ exists a plane $E \subset \RR^3$ with $A,B,C \in E$.
  For $A,B,C \in \RR^3$ with $A \not\in \geomline{BC}$ exists exactly one plane $E \subset \RR^3$ with $A,B,C \in E$.
\end{definition}
%____________________________________________________________________________________________________

\begin{theorem}[3-3.]\label{thm:3-3.}
  \uses{def:Postulate_1, def:Postulate_3, def:Collinear, def:Postulate_7}
  Given a line and a point not on the line, there is exactly one plane containing both of them.
  (\cite{SMSGGeometry}, 57)

  Given line $L \subset \RR^3$ and Point $P\in \RR^3$ with $P \not\in L$, there is exactly one Plane $E \subset \RR^3$ with $P \in E$ and $L \subset E$.
\end{theorem}

\begin{proof}
  \uses{def:Postulate_1, def:Postulate_3, def:Collinear, def:Postulate_7, def:Postulate_6}
  There exist $A,B \in L$ with $A\ne B$ (\ref{def:Postulate_3}) and $L=\geomline{AB}$.
  Since $P \not\in \geomline{AB}$ we know $A,B,C$ are non-collinear.
  From Postulate 7 (\ref{def:Postulate_7}) follows that there is exactly one plane containing the three points.
  From Postulate and (\ref{def:Postulate_6}) follows the line is also contained in that plane.

  %TODO proof seems hella wonky
\end{proof}

\begin{theorem}[3-4.]\label{thm:3-4.}
  \uses{def:Postulate_3, thm:3-1., def:Postulate_6, thm:3-3.}
  Given two intersecting lines, there is exactly one plane containing them.
  (\cite{SMSGGeometry}, 58)

  Let $L_1, L_2\subset \RR^3$ be two lines with $\#(L_1\cap L_2)=1$, then there is exactly one plane $E\subset \RR^3$ with $L_1,L_2 \subset E$.
\end{theorem}

\begin{proof}
  \uses{def:Postulate_3, thm:3-1., def:Postulate_6, thm:3-3., def:Collinear}
  Let $P \in L_1\cap L_2$ be the intersection.
  Py Postulate 3 (\ref{def:Postulate_3}) exists $Q \in L_1$ with $P \ne Q$.
  According to \ref{thm:3-1.} $Q \not\in L_2$.
  According to \ref{thm:3-3.} there is exactly one plane $E \subset \RR^3$ with $L_2 \subset E$ and $Q \in E$.
  By Postulate 6 (\ref{def:Postulate_6}), since $P,Q \in E$ follows $L_1=\geomline{PQ} \subset E$.
  Since $L_1, L_2$ define three non-collinear points $E$ is unique according to Postulate 7 (\ref{def:Postulate_7}).
\end{proof}

%_________________________________________ Postulate 8 _____________________________________________________________________
\begin{definition}[Postulate 8]\label{def:Postulate_8}
  \uses{def:GeomItems}
  If two different planes intersect then their intersection is a line.
  (\cite{SMSGGeometry}, 58)

  %TODO restatement possibly think about introducing dimensions?
\end{definition}
%________________________________________________________________________________________________________________________

\begin{definition}[Convex Sets]\label{def:Convex_Sets}
  \uses{def:Segment}
  A set $M \subset \RR^3$ is called convex if for every two points $P,Q \in A$ the entire segment $\segment{PQ}$ is in $M$.
  (\cite{SMSGGeometry}, 62)

  $M\subset \RR^3$ convex if $P,Q \in M \imp \segment{PQ}\subset A$.
\end{definition}


%_______________________________________________ Postulate 9 ______________________________________________________________
\begin{definition}[Postulate 9 - The Plane Separation Postulate]\label{def:Postulate_9}
  \uses{def:Segment, def:Postulate_6, def:Convex_Sets}
  Given a line and a plane containing it.
  The points of the plane that do not lie on the line form two sets such that
  \begin{itemize}
    \item each of the sets is convex and
    \item if $P$ is in one set and $Q$ in the other then $\segment{PQ}$ intersects the line.
  \end{itemize}
  (\cite{SMSGGeometry}, 64)

  Given line $L\subset\RR^3$ and plane $E\subset \RR^3$ with $L \subset E$.
  There exist two sets $M,N \subset\RR^3$ such that
  \begin{itemize}
    \item $M,N$ are convex and
    \item $P \in M \land Q \in N \imp \segment{PQ}\cap L \ne \emptyset$.
  \end{itemize}

  %TODO sets disjunkt?
\end{definition}
%___________________________________________________________________________________________________________________________

\begin{definition}[Half-Planes]\label{def:Half-Planes}
  \uses{def:Postulate_9}
  Given a line $L$ and a plane $E$ containing it, the two sets determined by Postulate 9 (\ref{def:Postulate_9}) are called half-planes, and $L$ is called an edge of each of them.
  We say that $L$ separates $E$ into the two half-planes.
  If two points lie in the same half plane, they are on the same side. Analogous: on opposite sides.
  (\cite{SMSGGeometry}, 64)

  Given a plane $E \subset \RR^3$ and two points $P,Q \in E$ and a line $L \in E$ we write
  $P\sim_L Q$ if $P$ and $Q$ are on the same side of $L$ and
  $P\not\sim_L Q$ if $P$ and $Q$ are on opposite sides of $L$.
  %TODO restatement?
\end{definition}

%_____________________________________ Postulate 10 _______________________________________________________________________
\begin{definition}[Postulate 10 - The Space Separation Postulate]\label{def:Postulate_10}
  \uses{def:GeomItems, def:Segment, def:Convex_Sets}
  The points of space that do not lie in a given plane form two sets such that
  \begin{itemize}
    \item each of the sets is convex and
    \item if a point $P$ is in one set and $Q$ in the other then $\segment{PQ}$ intersects the plane.
  \end{itemize}
  (\cite{SMSGGeometry}, 66)

  Given a plane $E \subset \RR^3$, there exist two sets $M,N \subset \RR^3$ such that
  \begin{itemize}
    \item $M,N$ are convex and
    \item $P\in M \land Q\in N \imp \segment{PQ}\cap E \ne \emptyset$.
  \end{itemize}
\end{definition}
%___________________________________________________________________________________________________________

\begin{definition}[Half-Spaces]\label{def:Half-Spaces}
  \uses{def:Postulate_10}
  The two sets determined by Postulate 10 (\ref{def:Postulate_10}) are called half-spaces, and the given plane is called the face of each of them.
  (\cite{SMSGGeometry}, 66)
  %TODO restatement?
\end{definition}

\begin{definition}[Angle]\label{def:Angle}
  \uses{def:Ray}
  An angle is the union of two rays (\ref{def:Ray}) which have the same end-point but do not lie in the same line.
  The two rays are called the sides of the angle and their common end-point is calles the vertex.
  (\cite{SMSGGeometry}, 71)

  For $\ray{AB}, \ray{AC} \subset \RR^3$ we define $\angle BAC:= \ray{AB} \cup \ray{AC}$.
\end{definition}

\begin{definition}[Triangle]\label{def:Triangle}
  \uses{def:Segment, def:Angle, def:Collinear}
  If $A,B,C$ are three non-collinear (\ref{def:Collinear}) points, then the union of the segments (\ref{def:Segment})
  $\segment{AB}, \segment{BC}, \segment{AC}$ is called a triangle, and is denoted $\geomtriang{ABC}$.
  The points $A,B,C$ are called its vertices, the segments its sides.
  $\angle ABC, \angle BAC, \angle ACB$ are called the angles of $\geomtriang{ABC}$.
  (\cite{SMSGGeometry}, 72)

  For $A,B,C \in \RR$ we define $\geomtriang{ABC}:=\segment{AB}\cup \segment{BC}\cup\segment{AC}$.
\end{definition}

\begin{definition}[Angle Interior]\label{def:Angle_Interior}
  \uses{def:Postulate_9, def:Half-Planes, def:Angle}
  Let $\angle BAC$ be an angle lying in a plane $E$.
  A point $P$ of $E$ lies in the interior of $\angle{BAC}$ if
  (1) $P$ and $B$ are on the same side of $\geomline{AC}$ and
  (2) $P$ and $C$ are on the same side of $\geomline{AB}$.
  The exterior of $\angle{BAC}$ is the set of all points that do not lie in the interior and do not lie on the angle itself.
  (\cite{SMSGGeometry}, 73)

  For $\angle{BAC} \subset $E$ \subset \RR^3$, with $E$ plane, we define
  \[
  \interior{\angle{BAC}}:=\set{P\mid P \in E \land P \sim_{\geomline{AC}} B \land P \sim_{\geomline{AB}} C}
  \]
  and
  \[
  \exterior{\angle{BAC}}:=E\setminus\left(\interior{\angle{BAC}}\cup \angle{BAC}\right) \MP.
  \]
\end{definition}

\begin{definition}[Triangle Interior]\label{def:Tiangle_Interior}
  \uses{def:Postulate_9, def:Half-Planes, def:Angle, def:Triangle, def:Angle_Interior}
  A point lies in the Interior of a triangle, if it lies in the interior of each of the angles of the triangle.
  A point lies in the exterior of the triangle if it lies in the plane of the triangle but is not a point of the triangle or its interior.
  (\cite{SMSGGeometry}, 74)

  Let $E \subset \RR^3$ be a plane and $\geomtriang{ABC}\subset E$ be a triangle.
  We define
  \[
  \interior{\geomtriang{ABC}}:=\interior{\angle{ABC}}\cap \interior{\angle{BAC}}\cap \interior{\angle{ACB}}
  \]
  and
  \[
  \exterior{\geomtriang{ABC}}:=E\setminus\left(\interior{\geomtriang{ABC}}\cup \geomtriang{ABC}\right)\MP.
  \]
\end{definition}

%_____________________________________________________ Postulate 11 __________________________________________________________________________
\begin{definition}[Postulate 11 - The Angle Measurement Postulate]\label{def:Postulate_11}
  \uses{def:Angle}
  To every angle $\angle{BAC}$ there corresponds a real number between 0 and 180.
  (\cite{SMSGGeometry}, 80)

  Let $\operatorname{ANG}\subset \RR^3$ be the set of all angles.
  There exists a function $\operatorname{m}:\operatorname{ANG}\to (0,180)$.
\end{definition}
%___________________________________________________________________________________________________________________________________________

\begin{definition}[Angle Measure]\label{def:Angle_Measure}
  \uses{def:Angle, def:Postulate_11}
  The number specified by Postulate 11 (\ref{def:Postulate_11}) is called the measure of the angle and written as $\measuredangle{ABC}$.
  (\cite{SMSGGeometry}, 80)

  For an angle $\angle{ABC} \subset \RR^3$ we define $\measuredangle{ABC}:=\operatorname{m}(\angle{ABC})$.
\end{definition}

%_____________________________________________________ Postulate 12 _______________________________________________________________________
\begin{definition}[Postulate 12 - The Angle Construction Postulate]\label{def:Postulate_12}
  \uses{def:Ray, def:Half-Planes, def:Angle_Measure}
  Let $\ray{AB}$ be a ray on the edge of the half-plane $H$.
  For every number $r$ between 0 and 180 there is exactly one ray $\ray{AP}$ with $P$ in $H$ such that $\measuredangle{PAB}=r$.
  (\cite{SMSGGeometry}, 81)

  Let $E \subset \RR^3$ be a plane, $\ray{AB}\in E$ be a ray and $Q,R \in E$ with $Q\not\sim_{\geomline{AB}}R$.
  For every $r \in (0,180)$ exists exactly one $P \in E$ with $P\sim_{\geomline{AB}}Q$ such that $\measuredangle{PAB}=r$.
\end{definition}
%_________________________________________________________________________________________________________________________________________



%_____________________________________________________ Postulate 13 _______________________________________________________________________
\begin{definition}[Postulate 13 - The Angle Addition Postulate]\label{def:Postulate_13}
  \uses{def:Angle, def:Angle_Interior, def:Angle_Measure}
  If $D$ is a point in the interior of $\angle{BAC}$, then
  \[
  \measuredangle{BAC}=\measuredangle{BAD}+\measuredangle{DAC}\MP.
  \]
  (\cite{SMSGGeometry}, 81)

  Let $\angle{BAC}\subset\RR^3$ be an angle.
  If $D \in \interior{\angle{BAC}}$ then $\measuredangle{BAC}=\measuredangle{BAD}+\measuredangle{DAC}$.
\end{definition}
%_________________________________________________________________________________________________________________________________________

\begin{definition}[Linear Pair]\label{def:Linear_Pair}
  \uses{def:Ray, def:Angle, def:Oppo_Rays}
  If $\ray{AB}$ and $\ray{AC}$ are opposite rays, and $\ray{AD}$ is another ray, than $\angle{BAD}$ and $\angle{DAC}$ form a linear pair.
  (\cite{SMSGGeometry}, 82)

  %TODO restatement
\end{definition}

\begin{definition}[Supplementary]\label{def:Supplementary}
  \uses{def:Angle_Measure, def:Angle}
  If the sum of the measure of two angles is 180 then the angles are called supplementary, and each is called supplement of the other.
  (\cite{SMSGGeometry}, 82)

  Let $\angle BAC, \angle EDF \subset \RR^3$ be angles. We define
  \[
  \supp{\angle BAC}{\angle DEF}:\eqv \measuredangle BAC + \measuredangle DEF=180\MP.
  \]
\end{definition}


%_____________________________________________________ Postulate 14 _______________________________________________________________________
\begin{definition}[Postulate 14 - The Supplement Postulate]\label{def:Postulate_14}
  \uses{def:Angle, def:Linear_Pair, def:Supplementary}
  If two angles form a linear pair, then they are supplementary.
  (\cite{SMSGGeometry}, 82)

  %TODO restatement
\end{definition}
%____________________________________________________________________________________________________________________________________________

\begin{definition}[Right Angle]\label{def:Right_Angle}
  \uses{def:Angle, def:Angle_Measure, def:Linear_Pair, def:Postulate_14}
  If the two angles of a linear pair have the same measure, then each of the angles is a right angle.
  By Postulate 14 (\ref{def:Postulate_14}) follows that a right angle is an angle $\angle A$ with $\measuredangle A=90$.
  (\cite{SMSGGeometry}, 85)

  %TODO symbol for right angle?
\end{definition}

\begin{definition}[Perpendicular]\label{def:Perpendicular}
  \uses{def:Right_Angle} %TODO also segment, ray definition? line that contains segment or ray depends on Postulate 1?
  Two intersecting sets, each of which is either a line, a ray or a segment, are perpendicular if the two lines which contain them determine a right angle.
  (\cite{SMSGGeometry}, 86)

  %TODO restatement... how to denote one dimensional subsets?
  We write $\perp$.

\end{definition}

\begin{definition}[Complementary]\label{def:Complementary}
  \uses{def:Angle_Measure}
  If the sum of the measures of two angles is 90, then the angles are called complementary, and each of them is the complement of the other.
  (\cite{SMSGGeometry}, 86)
  Let $\angle BAC, \angle EDF \subset \RR^3$ be angles. We define
  \[
  \comp{\angle BAC}{\angle DEF}:\eqv \measuredangle BAC + \measuredangle DEF=90\MP.
  \]

\end{definition}

\begin{definition}[Acute-Obtuse]\label{def:Acute_Obtuse}
  \uses{def:Angle_Measure}
  An angle with measure less than 90 is called acute.
  An angle with measure greater than 90 is called obtuse.
  (\cite{SMSGGeometry}, 86)

  %TODO restatement necessary?
\end{definition}

\begin{definition}[Congruent Angles]\label{def:Congruent_Angles}
  \uses{def:Angle_Measure}
  Angles with the same measure are called congruent.
  (\cite{SMSGGeometry}, 86)

  For two angles $\angle BAC, \angle EDF$ let
  \[
  \angle BAC \cong \angle EDF :\eqv \measuredangle BAC = \measuredangle EDF\MP.
  \]

\end{definition}

\begin{theorem}[4-1.]\label{thm:4-1.}
  \uses{def:Angle, def:Angle_Measure, def:Complementary, def:Acute_Obtuse}
  If two angles are complementary, then both of them are acute.
  (\cite{SMSGGeometry}, 87)

  For $\angle BAC, \angle EDF \subset \RR^3: \comp{\angle ABC}{\angle EDF} \imp \measuredangle BAC, \measuredangle EDF <90$.
\end{theorem}

\begin{proof}
  \uses{def:Angle, def:Angle_Measure, def:Complementary, def:Acute_Obtuse}
  Let $\angle BAC, \angle EDF \subset \RR^3$ with $ \comp{\angle ABC}{\angle EDF}$:
  \begin{align*}
    \comp{\angle ABC}{\angle EDF}
        &\impu{\text{\ref{def:Complementary}}} \measuredangle BAC+\measuredangle EDF = 90 \\
        &\imp \measuredangle BAC=90-\measuredangle EDF \land \measuredangle EDF=90-\measuredangle BAC \\
        &\impu{\text{\ref{def:Postulate_11}}} \measuredangle BAC <90 \land \measuredangle EDF < 90 \\
        &\impu{\text{\ref{def:Acute_Obtuse}}} \angle BAC, \angle EDF \text{ acute.}
  \end{align*}
  %TODO make sure ref works in align environment
\end{proof}

\begin{theorem}[4-2.]\label{thm:4-2.}
  \uses{def:Angle, def:Postulate_11, def:Congruent_Angles}
  Every angle is congruent to itself.
  (\cite{SMSGGeometry}, 87)

  For all $\angle BAC \subset \RR^3: \angle BAC \cong \angle BAC$.
\end{theorem}

\begin{proof}
  \uses{def:Angle, def:Postulate_11, def:Congruent_Angles}
  Postulate 11 (\ref{def:Postulate_11}) defines function $\operatorname{m}$, meaning $\operatorname{m}(\angle BAC)=\measuredangle BAC$ is unique.
  Hence $\measuredangle BAC = \measuredangle BAC \imp \angle BAC \cong \angle BAC$.
  %TODO proof seems wonky
\end{proof}

\begin{theorem}[4-3.]\label{thm:4-3.}
  \uses{def:Angle_Measure, def:Right_Angle, def:Congruent_Angles}
  Any two right angles are congruent.
  (\cite{SMSGGeometry}, 87)

  For any angles $\angle A, \angle B \subset \RR^3$ with $\measuredangle A=\measuredangle B=90$ follows $\angle A \cong \angle B$.
\end{theorem}

\begin{proof}
  \uses{def:Angle_Measure, def:Right_Angle, def:Congruent_Angles}
  Let $\angle A, \angle B \subset \RR^3$ be angles with $\measuredangle A=\measuredangle B = 90$.
  From $\measuredangle A = \measuredangle B$ follows from (\ref{def:Congruent_Angles}) that $\angle A \cong \angle B$.
\end{proof}

\begin{theorem}[4-4.]\label{thm:4-4.}
  \uses{def:Angle, def:Angle_Measure, def:Supplementary, def:Congruent_Angles, def:Right_Angle}
  If two angles are congruent and supplementary, then each of them is a right angle.
  (\cite{SMSGGeometry}, 87)

  Let $\angle A, \angle B\subset \RR^3$ be angles.
  If $\angle A \cong \angle B \land \supp{\angle A}{\angle B}$ then $\angle A= \angle B =90$.
\end{theorem}

\begin{proof}
  \uses{def:Angle, def:Angle_Measure, def:Supplementary, def:Congruent_Angles, def:Right_Angle}
  \begin{align}
    \angle A \cong \angle B
        &\imp \measuredangle A = \measuredangle B =:r \MP, \label{4-4.1}\\
    \supp{\angle A}{\angle B}
        &\imp \measuredangle A + \measuredangle B = 180 \MP, \label{4-4.2}\\
    \ref{4-4.1},\ref{4-4.2}
        &\imp \measuredangle A + \measuredangle B = 2r =180 \\
        &\imp \measuredangle A + \measuredangle B = r=90 \MP. \\
      \end{align}
  %TODO sort out label and ref, it's not working
\end{proof}

\begin{theorem}[4-5.]\label{thm:4-5.}
  \uses{def:Angle_Measure, def:Supplementary, def:Congruent_Angles}
  Supplements of congruent angles are congruent.
  (\cite{SMSGGeometry}, 87)

  Let $\angle A,\angle B, \angle C \subset \RR^3$ be angles.
  If $\angle B \cong \angle C, \supp{\angle A}{\angle B}$ and $\supp{\angle C}{\angle D}$ then $\angle A \cong \angle C$.
\end{theorem}

\begin{proof}
  \uses{def:Angle_Measure, def:Supplementary, def:Congruent_Angles}
  \begin{align*}
    \angle B \cong \angle D &\imp \measuredangle B=\measuredangle D =:r\MP,\\
    \supp{\angle A}{\angle B}&\imp \measuredangle A=180-\measuredangle B= 180 -r \MP,\\
    \supp{\angle C}{\angle D} &\imp \measuredangle C = 180 -\measuredangle D=180-r \MP,\\
    \measuredangle A=\measuredangle C=180-r &\imp \angle A\cong \angle C\MP.
  \end{align*}
\end{proof}

\begin{theorem}[4-6.]\label{thm:4-6.}
  \uses{def:Angle_Measure, def:Complementary, def:Congruent_Angles}
  Complements of congruent angles are congruent.
  (\cite{SMSGGeometry}, 88)

  Let $\angle A,\angle B, \angle C \subset \RR^3$ be angles.
  If $\angle B \cong \angle C, \comp{\angle A}{\angle B}$ and $\comp{\angle C}{\angle D}$ then $\angle A \cong \angle C$.
\end{theorem}

\begin{proof}
  Analogous to proof for \ref{thm:4-5.}.
\end{proof}

\begin{definition}[Vertical Angles]\label{def:Vertical_Angles}
  \uses{def:Angle, def:Oppo_Rays}
  Two angles are vertical angles, if their sides form two pairs of opposite rays.
  (\cite{SMSGGeometry}, 88)

  Write $\vertAng{\angle A}{\angle B}$.

  %TODO restatement... might rely on restatement of Oppo_Rays, might have to find symbol for opposite rays to make things easier
\end{definition}

\begin{theorem}[4-7.]\label{thm:4-7.}
  \uses{def:Oppo_Rays, def:Angle, def:Supplementary, def:Postulate_14, def:Linear_Pair, def:Congruent_Angles}
  Vertical Angles are congruent.
  (\cite{SMSGGeometry}, 88)

  For angles $\angle A,\angle B \subset \RR^3$ with $\vertAng{\angle A}{\angle B}$ follows $\angle A \cong \angle B$.
\end{theorem}

\begin{proof}
  \uses{def:Oppo_Rays, def:Angle, def:Supplementary, def:Postulate_14, def:Linear_Pair, def:Congruent_Angles}
  %TODO write proof
\end{proof}

\begin{theorem}[4-8.]\label{thm:4-8.}
  \uses{def:Postulate_3, def:Oppo_Rays, def:Angle_Measure, def:Linear_Pair, def:Supplementary, def:Postulate_14, def:Right_Angle, def:Vertical_Angles, thm:4-7.} %TODO uses def:Postulate_13?
  If two intersecting lines form one right angle, then they form four right angles.
  %TODO restatement
\end{theorem}

\begin{proof}
  %TODO
\end{proof}












\section{Bibliography}
	\bibliographystyle{plain}
	\begin{thebibliography}{1}

	\bibitem{SMSGGeometry}
    School Mathematics Study Group, \textit{Geometry: Unit 13 -- Student's Text, Part I}, 1961


	\end{thebibliography}
